\section{Reasoning scaffold：让中间过程可搜索、可检查}

课程最后一段模型内容讨论 Chain-of-Thought、Tree of Thoughts 与 Socratic Questioning。共同思路是把“一次直接输出”改成多步中间过程，使模型获得更多 test-time compute，也给系统留下检查和回退的位置。

\subsection{CoT 为什么可能有用}

Chain-of-Thought（思维链，CoT）通过示例或指令要求模型生成中间步骤。对多步算术和符号任务，直接映射问题到答案很难；把任务分解后，每一步局部关系可能更接近训练分布。接下来的两页先展示 intermediate guidance，再比较 standard prompt 与 CoT example 的信息差。

\lecturefigure{slide-070.jpg}{Intermediate guidance 将一次预测展开为多步推理}{官方 CS25 V4 slide deck，第 70 页}

\lecturefigure{slide-071.jpg}{标准 prompting 与 Chain-of-Thought prompting 对比}{官方 CS25 V4 slide deck，第 71 页}

\begin{importantbox}{CoT 的概率分解}
令中间 reasoning trace 为 $r$、最终答案为 $y$，生成过程可写成
\[
p(y\mid x)=\sum_r p(r\mid x)p(y\mid x,r).
\]
实际 decoding 只采样或搜索少量 $r$。CoT 的优势来自更丰富的中间计算与分解，不保证采样到的 trace 真实反映模型内部因果过程。
\end{importantbox}

\teachervoice{讲者用十位数乘法说明中间步骤的直觉：人类通常也要分步计算。CoT 的额外价值是错误可见，读者能定位模型在哪一步遗漏或误解，而不是只看到一个错误答案。}

\subsection{性能提升之外，要看错误类型}

CoT 在大模型上常有提升，但仍会生成无关步骤、局部错误或正确过程后的错误结论。课堂把错误分成 missing step 与 semantic misunderstanding，尤其强调小模型可能连局部 mapping 都不稳定。本节读图重点是 error taxonomy 与 model size，而不是只记住平均性能上升。

\lecturefigure{slide-072.jpg}{CoT 提升与剩余错误的课堂总结}{官方 CS25 V4 slide deck，第 72 页}

\begin{warningbox}{长 rationale 不等于可靠 reasoning}
模型可以生成格式完整却与答案无关的解释，也可以在得到答案后反向编写理由。验收应看中间状态是否可执行、可验证、能否支持 counterfactual change，而不是只看文字是否像推理。
\end{warningbox}

\lecturefigure{slide-073.jpg}{小模型的 missing-step 与 semantic-understanding 错误}{官方 CS25 V4 slide deck，第 73 页}

\begin{knowledgebox}{为什么小模型更容易在 CoT 中失败}
CoT 增加了序列长度和必须保持的一致性。若模型缺少基本语义或算术能力，更多 token 会创造更多出错位置。改进路线包括分步 supervision、tool execution、verifier、distillation 与更结构化的 state representation。
\end{knowledgebox}

\subsection{Generalized CoT：推理不只有一条线}

课程提出 CoT 定义过于刚性：不同问题可能需要分支搜索、回溯、子问题递归或外部工具。Generalized reasoning 的重点是选择合适的 computation graph。下面三页从线性 trace 过渡到 tree search 和 recursive decomposition，比较 state、transition 与 verifier。

\lecturefigure{slide-074.jpg}{从固定 Chain-of-Thought 走向更一般的推理结构}{官方 CS25 V4 slide deck，第 74 页}

\begin{importantbox}{Reasoning scaffold 的三个设计轴}
\begin{enumerate}
  \item \textbf{State}：中间状态是自然语言、程序、方程还是环境 observation；
  \item \textbf{Transition}：下一步由模型采样、规则、搜索或工具执行产生；
  \item \textbf{Verifier}：局部状态怎样评分、剪枝、回退与终止。
\end{enumerate}
只有把三者写清楚，推理方法才不只是 prompt 风格。
\end{importantbox}

\subsection{Tree of Thoughts 与 Socratic decomposition}

Tree of Thoughts（思维树，ToT）维护多个候选状态并进行 lookahead/backtracking；Socratic Questioning（苏格拉底式提问）递归提出子问题，再把答案组合回原问题。两者都用更多 inference compute 换取搜索空间。

\lecturefigure{slide-075.jpg}{Tree of Thoughts 的分支、评估与回溯}{官方 CS25 V4 slide deck，第 75 页}

\begin{importantbox}{有限宽度搜索}
若每个状态扩展 $b$ 个候选、保留 top-$k$、深度为 $d$，粗略模型调用量随 $O(dkb)$ 增长。Verifier 的偏差会决定哪些分支被过早删除，因此搜索质量不仅取决于 generator，还取决于 state scoring。
\end{importantbox}

\lecturefigure{slide-076.jpg}{Socratic Questioning 递归生成并回答子问题}{官方 CS25 V4 slide deck，第 76 页}

\begin{knowledgebox}{读图：decomposition 与 search 的差别}
Decomposition 把大问题拆成相对独立的小问题；search 在多个候选路径中选择。Socratic method 更强调子问题结构，ToT 更强调候选状态和回溯。复杂系统常把两者结合，并在叶节点调用 calculator、code executor 或 retrieval。
\end{knowledgebox}

\subsection{本章小结}

Reasoning scaffold 的价值在于增加中间计算、暴露错误并提供搜索与验证接口。CoT、ToT 和 Socratic decomposition 并不自动产生真 reasoning；它们只有与结构化状态、tool、verifier 和 stopping rule 结合时，才更接近可控计算过程。

